\documentclass[11pt,a4paper]{article}

% ------------------------------------------------
% CORE LAYOUT & TYPOGRAPHY (publication grade)
% ------------------------------------------------
\usepackage[margin=1in]{geometry}
\usepackage[T1]{fontenc}
\usepackage[utf8]{inputenc}
\usepackage{lmodern}                          % high-quality Latin Modern
\usepackage[final]{microtype}                 % protrusion & expansion
\usepackage{setspace}
\onehalfspacing                               % comfortable line spacing
\usepackage{parskip}                          % no paragraph indent, vertical space
\usepackage{ragged2e}
\usepackage{enumitem}
\setlist{itemsep=0.35\baselineskip, topsep=0.35\baselineskip}
\usepackage{multicol}
\setlength{\columnsep}{18pt}
\usepackage{booktabs}
\usepackage{array}
\usepackage{tabularx}
\usepackage{makecell}
\usepackage{caption}
\usepackage{xcolor}
\usepackage{csquotes}
\usepackage{siunitx}
\usepackage{titlesec}
\usepackage{fancyhdr}
\usepackage{hypcap}
\usepackage{hyperref}
\usepackage{cleveref}

% ------------------------------------------------
% TIKZ FOR DIAGRAMS
% ------------------------------------------------
\usepackage{tikz}
\usetikzlibrary{shapes.geometric, arrows.meta, positioning, fit, backgrounds, calc, shadows}
\tikzset{
  block/.style={rectangle, draw, fill=blue!10, text width=6em, text centered, rounded corners, minimum height=3em, drop shadow},
  decision/.style={diamond, draw, fill=green!10, text width=4.5em, text badly centered, inner sep=0pt, drop shadow},
  line/.style={draw, -Latex, thick},
  cloud/.style={draw, ellipse, fill=red!10, minimum height=2em, drop shadow},
  process/.style={rectangle, draw, fill=orange!10, text width=8em, text centered, rounded corners, minimum height=3em, drop shadow},
  smallblock/.style={rectangle, draw, fill=blue!10, text width=4.5em, text centered, rounded corners, minimum height=2.5em, font=\small},
}

% ------------------------------------------------
% ALGORITHMS
% ------------------------------------------------
\usepackage{algorithm}
\usepackage{algpseudocode}
\algrenewcommand\algorithmicrequire{\textbf{Input:}}
\algrenewcommand\algorithmicensure{\textbf{Output:}}
\algrenewcommand\algorithmiccomment[1]{\hfill\(\triangleright\)~#1}

% ------------------------------------------------
% CODE: minted (best for publishing)
% ------------------------------------------------
\usepackage[cache=false]{minted}
\usepackage{fvextra}
\definecolor{codebg}{RGB}{248,248,248}
\definecolor{codeframe}{RGB}{225,225,225}
\setminted{
  fontsize=\small,
  baselinestretch=1.1,
  breaklines,
  breakanywhere,
  autogobble,
  bgcolor=codebg,
  frame=lines,
  framerule=0.4pt,
  framesep=6pt,
  rulecolor=codeframe,
  tabsize=2,
  encoding=utf8,
  obeytabs=false
}
% Language helpers
\newmintedfile{text}{}
\newminted{c}{}
\newminted{asm}{}
\newminted{text}{}

% ------------------------------------------------
% HYPERREF
% ------------------------------------------------
\hypersetup{
  colorlinks=true,
  linkcolor=blue!50!black,
  citecolor=blue!50!black,
  urlcolor=blue!60!black,
  pdfauthor={Md Sadek Hossain Asif, Tanzila Khan},
  pdftitle={DUOS Assignment 02 - System Call and Exception Handling Report},
  pdfsubject={Cortex-M4F SVC-based syscall ABI and exception path},
  pdfkeywords={DUOS, STM32F446RE, Cortex-M4F, SysTick, PendSV, SVC, system calls}
}
\urlstyle{same}

% ------------------------------------------------
% PAGE STYLE & TITLES
% ------------------------------------------------
\pagestyle{fancy}
\fancyhf{}
\lhead{\textit{DUOS Assignment 02}}
\rhead{\thepage}
\setlength{\headheight}{14pt}

\titleformat{\section}{\large\bfseries}{\thesection.}{0.7em}{}
\titleformat{\subsection}{\normalsize\bfseries}{\thesubsection.}{0.6em}{}
\titleformat{\subsubsection}{\normalsize\bfseries\itshape}{\thesubsubsection.}{0.6em}{}
\titlespacing*{\section}{0pt}{0.7\baselineskip}{0.45\baselineskip}
\titlespacing*{\subsection}{0pt}{0.55\baselineskip}{0.35\baselineskip}
\titlespacing*{\subsubsection}{0pt}{0.5\baselineskip}{0.3\baselineskip}

\captionsetup{labelfont=bf,font=small,tableposition=top}

% ------------------------------------------------
% TABLE HELPERS: flexible wrapped columns
% ------------------------------------------------
\newcolumntype{L}[1]{>{\RaggedRight\arraybackslash}p{#1}}
\newcolumntype{C}[1]{>{\Centering\arraybackslash}p{#1}}
\newcolumntype{R}[1]{>{\RaggedLeft\arraybackslash}p{#1}}
\newcolumntype{Y}{>{\RaggedRight\arraybackslash}X}

% ------------------------------------------------
% MACROS (kept from your original)
% ------------------------------------------------
\newcommand{\file}[1]{\texttt{#1}}
\newcommand{\reg}[1]{\texttt{#1}}

% ------------------------------------------------
% TOC DEPTH
% ------------------------------------------------
\setcounter{secnumdepth}{3}
\setcounter{tocdepth}{3}

% ------------------------------------------------
% SAFETY: avoid overfull/underfull + footnote splits
% ------------------------------------------------
\sloppy
\emergencystretch=3em
\clubpenalty=10000
\widowpenalty=10000
\interfootnotelinepenalty=10000
\frenchspacing

% ------------------------------------------------
% TITLE PAGE
% ------------------------------------------------
\begin{document}

% Custom title page
\begin{titlepage}
  \centering
  \vspace*{1.5cm}
  
  {\Large\textbf{University of Dhaka}}\\[0.3cm]
  {\large Department of Computer Science and Engineering}\\[2cm]
  
  \begin{tikzpicture}
    \draw[line width=1pt] (0,0) -- (12,0);
  \end{tikzpicture}\\[0.5cm]
  
  {\huge\bfseries DUOS Assignment 02}\\[0.4cm]
  {\LARGE System Call and Exception Handling Report}\\[0.5cm]
  
  \begin{tikzpicture}
    \draw[line width=1pt] (0,0) -- (12,0);
  \end{tikzpicture}\\[1cm]
  
  {\large\textit{Implementation of SVC-based System Call ABI}}\\
  {\large\textit{for STM32F446RE Cortex-M4F}}\\[2.5cm]
  
  \begin{minipage}{0.45\textwidth}
    \begin{flushleft}
      \textbf{\large Submitted By:}\\[0.3cm]
      Md Sadek Hossain Asif\\
      Tanzila Khan\\
      Department of CSE
    \end{flushleft}
  \end{minipage}
  \hfill
  \begin{minipage}{0.45\textwidth}
    \begin{flushright}
      \textbf{\large Submitted To:}\\[0.3cm]
      Mossaddek Kamal Tushar\\
      Jargis Ahmed\\
      Department of CSE
    \end{flushright}
  \end{minipage}
  
  \vfill
  
  {\large \today}
  
\end{titlepage}

\tableofcontents
\newpage



\subsection*{Objectives}
\begin{itemize}[leftmargin=1.5em]
  \item Introduce a software interrupt (SVC) driven system call ABI for DUOS on the STM32F446RE (Cortex--M4F).
  \item Provide the seven mandatory services (\texttt{read}, \texttt{write}, \texttt{time}, \texttt{getpid}, \texttt{yield}, \texttt{reboot}, \texttt{exit}) to userland in a controlled manner.
  \item Configure SysTick and PendSV to build the scheduling foundation for later assignments.
  \item Demonstrate the end-to-end request path using a userland demo application and provide memory usage statistics.
\end{itemize}



\newpage

\section{Introduction}

\subsection{Background}
Modern operating systems require a well-defined mechanism to transition between unprivileged user code and privileged kernel services. The ARM Cortex-M4F architecture provides the Supervisor Call (SVC) exception specifically for this purpose, enabling user applications to request kernel services safely without compromising system integrity.

\subsection{System Architecture Overview}

\begin{figure}[h]
\centering
\begin{tikzpicture}[node distance=1.5cm]
  % User Space
  \node[process] (user) {User Application};
  \node[process, below=of user] (wrapper) {Syscall Wrappers\\(unistd.c)};
  
  % Hardware boundary
  \draw[thick, dashed] ([xshift=-4cm, yshift=-1cm]wrapper.south) -- ([xshift=4cm, yshift=-1cm]wrapper.south) node[midway, above, fill=white] {Hardware Transition (SVC)};
  
  % Kernel Space
  \node[block, below=2cm of wrapper] (svhandler) {SVCall\\Handler};
  \node[process, below=of svhandler] (decoder) {SVC\_Handler\_C\\(Decoder)};
  \node[process, below=of decoder] (dispatcher) {syscall\_dispatch()};
  
  % Services
  \node[smallblock, below left=1.2cm and 1.5cm of dispatcher] (uart) {UART\\Driver};
  \node[smallblock, below=1.2cm of dispatcher] (systick) {SysTick\\Timer};
  \node[smallblock, below right=1.2cm and 1.5cm of dispatcher] (system) {System\\Control};
  
  % Arrows
  \draw[line] (user) -- (wrapper);
  \draw[line] (wrapper) -- node[right] {\texttt{svc \#imm}} (svhandler);
  \draw[line] (svhandler) -- (decoder);
  \draw[line] (decoder) -- (dispatcher);
  \draw[line] (dispatcher) -- (uart);
  \draw[line] (dispatcher) -- (systick);
  \draw[line] (dispatcher) -- (system);
  
  % Return path
  \draw[line, dashed, blue] (dispatcher.east) -- ++(1.5,0) |- node[near start, right] {Return} (svhandler.east);
  
  % Labels
  \node[above left=0.2cm and 0.5cm of user, text=blue!70!black] {\textbf{User Space}};
  \node[below left=0.2cm and 0.5cm of uart, text=red!70!black] {\textbf{Kernel Space}};
\end{tikzpicture}
\caption{DUOS System Call Architecture Overview}
\label{fig:architecture}
\end{figure}

\subsection{Design Goals}
\begin{enumerate}[leftmargin=2em]
  \item \textbf{Safety:} Enforce privilege separation through hardware exception mechanisms.
  \item \textbf{Efficiency:} Minimize context switching overhead and maintain deterministic timing.
  \item \textbf{Compatibility:} Conform to ARM EABI calling conventions for future extensibility.
  \item \textbf{Clarity:} Maintain clean layering between user wrappers, exception handlers, and kernel services.
\end{enumerate}

\section{System Call Flow}

\subsection{High-Level Flow}
Figure~\ref{fig:detailed-flow} and \Cref{alg:svc-path} detail how a user request traverses the system. Each step corresponds to concrete code references in the subsequent sections.
\begin{figure}[H]
\centering
\begin{tikzpicture}[node distance=1.3cm, every node/.style={font=\small}]
  % Step 1
  \node[process, text width=7cm] (step1) {%
    \textbf{Step 1:} User calls wrapper\\
    \texttt{write(STDOUT\_FILENO, buf, len)}\\
    \textit{Location:} \file{userland/utils/unistd.c}
  };
  
  % Step 2
  \node[process, below=of step1, text width=7cm] (step2) {%
    \textbf{Step 2:} Load registers \& invoke SVC\\
    \texttt{r0-r3} $\leftarrow$ arguments\\
    \texttt{svc \#55} (SYS\_write)
  };
  
  % Step 3
  \node[block, below=of step2, text width=7cm] (step3) {%
    \textbf{Step 3:} Hardware exception entry\\
    CPU $\rightarrow$ Handler Mode (MSP)\\
    Auto-stack: r0-r3, r12, LR, PC, xPSR
  };
  
  % Step 4
  \node[process, below=of step3, text width=7cm] (step4) {%
    \textbf{Step 4:} SVCall\_Handler (ASM)\\
    Detect stack (MSP/PSP)\\
    Pass frame pointer to C handler
  };
  
  % Step 5
  \node[process, below=of step4, text width=7cm] (step5) {%
    \textbf{Step 5:} SVC\_Handler\_C\\
    Decode SVC immediate\\
    Extract arguments from stack
  };
  
  % Step 6
  \node[process, below=of step5, text width=7cm, fill=green!10] (step6) {%
    \textbf{Step 6:} syscall\_dispatch()\\
    Execute kernel service\\
    \texttt{case SYS\_write: ...}
  };
  
  % Step 7
  \node[process, below=of step6, text width=7cm] (step7) {%
    \textbf{Step 7:} Write return to stack\\
    \texttt{stack[0] = retval}\\
    Exception return (EXC\_RETURN)
  };
  
  % Step 8
  \node[process, below=of step7, text width=7cm] (step8) {%
    \textbf{Step 8:} Hardware exception exit\\
    Restore context from stack\\
    Return to Thread Mode with result in r0
  };
  
  % Arrows
  \draw[line] (step1) -- (step2);
  \draw[line] (step2) -- (step3);
  \draw[line] (step3) -- (step4);
  \draw[line] (step4) -- (step5);
  \draw[line] (step5) -- (step6);
  \draw[line] (step6) -- (step7);
  \draw[line] (step7) -- (step8);
  
  % Side annotations
  \node[right=0.3cm of step1, text width=3cm, align=left, text=blue] {User Space\\(Thread Mode)};
  \node[right=0.3cm of step3, text width=3cm, align=left, text=red] {Kernel Space\\(Handler Mode)};
\end{tikzpicture}
\caption{Detailed System Call Execution Flow}
\label{fig:detailed-flow}
\end{figure}

\subsection{Step-by-Step Description}
\begin{enumerate}[leftmargin=2em,label=\arabic*.]
  \item User code calls a libc-style wrapper (e.g. \texttt{write()}) in \file{userland/utils/unistd.c}.
  \item The wrapper populates \reg{r0--r3}, issues \texttt{svc \#imm}, and the core transitions to handler mode.
  \item \texttt{SVCall\_Handler} in \file{kern/arch/stm32f446re/sys\_lib/stm32\_startup.c} selects the correct stack (MSP/PSP) and branches to the C-level helper with the hardware-stacked frame pointer.
  \item \texttt{SVC\_Handler\_C} decodes the immediate operand, gathers the arguments, and forwards the call to \texttt{syscall\_dispatch()} in \file{kern/syscall/syscall.c}.
  \item The dispatcher executes the requested kernel service. UART access is funnelled through \file{kern/arch/stm32f446re/sys\_lib/sys\_usart.c}, timing relies on \file{kern/arch/cm4/systick.c}, and termination/yield manipulates \reg{SCB->ICSR}.
  \item The dispatcher return value is written back into the stacked \reg{r0}. When the exception unwinds, the user thread sees the correct result and resumes execution in thread mode.
\end{enumerate}

\subsection{Pseudocode Algorithm}
\begin{algorithm}[h]
  \caption{End-to-End SVC Request Path}\label{alg:svc-path}
  \begin{algorithmic}[1]
    \Require User wrapper (e.g., \texttt{write()}) with args in \(\texttt{r0..r3}\)
    \Ensure Result propagated back to user \(\texttt{r0}\)
    \State \textbf{Userland}: Call wrapper in \file{userland/utils/unistd.c}
    \State Wrapper issues \texttt{svc \#imm} \Comment{Core enters Handler mode}
    \State \textbf{SVCall\_Handler} (\file{stm32\_startup.c}): select MSP/PSP, pass stacked frame to C
    \State \textbf{SVC\_Handler\_C}: decode immediate, extract \(\texttt{a0..a3}\), call \texttt{syscall\_dispatch()}
    \State \textbf{Dispatcher} (\file{kern/syscall/syscall.c}): execute kernel service
    \State Write return value to stacked \(\texttt{r0}\); unwind exception to Thread mode
    \State \Return to user with correct result in \(\texttt{r0}\)
  \end{algorithmic}
\end{algorithm}

\subsection{Exception Stack Frame}
Figure~\ref{fig:stack-frame} illustrates the stack layout during SVC exception handling.

\begin{figure}[H]
\centering
\begin{tikzpicture}[
  stackframe/.style={rectangle, draw, fill=yellow!20, minimum width=4cm, minimum height=0.7cm, font=\small},
  stacklabel/.style={font=\small\ttfamily}
]
  % Stack frames
  \node[stackframe] (xpsr) at (0,0) {xPSR};
  \node[stackframe] (pc) at (0,-0.7) {PC (Return Address)};
  \node[stackframe] (lr) at (0,-1.4) {LR};
  \node[stackframe] (r12) at (0,-2.1) {R12};
  \node[stackframe] (r3) at (0,-2.8) {R3 (arg3)};
  \node[stackframe] (r2) at (0,-3.5) {R2 (arg2)};
  \node[stackframe] (r1) at (0,-4.2) {R1 (arg1)};
  \node[stackframe] (r0) at (0,-4.9) {R0 (arg0/retval)};
  
  % Labels
  \node[left=0.5cm of xpsr, stacklabel] {[7]};
  \node[left=0.5cm of pc, stacklabel] {[6]};
  \node[left=0.5cm of lr, stacklabel] {[5]};
  \node[left=0.5cm of r12, stacklabel] {[4]};
  \node[left=0.5cm of r3, stacklabel] {[3]};
  \node[left=0.5cm of r2, stacklabel] {[2]};
  \node[left=0.5cm of r1, stacklabel] {[1]};
  \node[left=0.5cm of r0, stacklabel] {[0]};
  
  % Stack pointer
  \draw[-Latex, thick, red] (3, -4.9) -- (2, -4.9) node[midway, above] {SP};
  
  % Growth direction
  \draw[-Latex, thick, blue] (5, 0.5) -- (5, -5.4) node[midway, right] {Stack Growth};
  
  % Annotations
  \node[right=0.5cm of pc, text width=4cm, align=left, font=\small] {Auto-stacked by\\hardware on\\exception entry};
  \node[right=0.5cm of r0, text width=4cm, align=left, font=\small, text=red] {Modified by\\kernel for\\return value};
\end{tikzpicture}
\caption{Exception Stack Frame Layout (ARMv7-M)}
\label{fig:stack-frame}
\end{figure}

\section{Boot and Core Services}

\subsection{System Initialisation}
\file{kern/lib/kern/sys\_init.c} (\texttt{\_\_sys\_init()}) performs the mandatory hardware bring-up:
\begin{itemize}[leftmargin=1.5em]
  \item 180\,MHz system clock, FPU enable, nested interrupt priority grouping.
  \item 1\,kHz SysTick via \texttt{\_\_SysTick\_init(1000)}.
  \item USART2 (console) and USART6, ring buffer initialisation, timer configuration.
  \item Boot banner and system diagnostics in debug builds.
\end{itemize}
The demo's \texttt{kmain()} (\file{kern/kmain/kmain.c}) calls \texttt{\_\_sys\_init()}, seeds the global task identifier, and orchestrates syscall probes once userland wrappers become available.

\subsection{SysTick Handler}
\file{kern/arch/cm4/systick.c} maintains the monotonic millisecond counter used by \texttt{SYS\_\_\_time}:

\begin{figure}[H]
\centering
\begin{tikzpicture}[node distance=1.2cm]
  \node[cloud] (systick) {SysTick IRQ\\(1 kHz)};
  \node[process, below=of systick, text width=6cm] (increment) {Increment \texttt{g\_msTick}};
  \node[decision, below=of increment, aspect=2] (check) {Every\\10 ms?};
  \node[block, below=of check, text width=5cm] (pendsv) {Set PendSV\\SCB->ICSR |= PENDSVSET};
  \node[block, right=2cm of check] (continue) {Continue};
  
  \draw[line] (systick) -- (increment);
  \draw[line] (increment) -- (check);
  \draw[line] (check) -- node[left] {Yes} (pendsv);
  \draw[line] (check) -- node[above] {No} (continue);
  \draw[line] (pendsv) -| (continue);
\end{tikzpicture}
\caption{SysTick and PendSV Triggering Logic}
\label{fig:systick-flow}
\end{figure}

\begin{itemize}[leftmargin=1.5em]
  \item \texttt{g\_msTick} increments every interrupt, providing \texttt{\_\_getTime()} and utility accessors (seconds, minutes, hours).
  \item A secondary counter (\texttt{tick10}) triggers PendSV every 10\,ms via \reg{SCB->ICSR |= PENDSVSET}, cementing the scheduling cadence for the forthcoming context switcher.
  \item The interrupt priority is explicitly lowered (0xFF) to keep SysTick and PendSV preemptible by real-time ISRs.
\end{itemize}

\section{SVC Entry Path}

\subsection{Vector Table Updates}
The device vector table in \file{kern/arch/stm32f446re/sys\_lib/stm32\_startup.c} now points the SVCall slot at a custom naked handler and connects PendSV to a stub that currently returns immediately (``bx lr''). This preserves a clean hook for scheduler logic while avoiding the default weak implementations.

\subsection{Assembly Shim}
The SVCall\_Handler is implemented as a naked function in assembly to efficiently detect which stack pointer was in use and pass it to the C handler.

\vspace{0.3cm}
\noindent\textbf{File:} \file{kern/arch/stm32f446re/sys\_lib/stm32\_startup.c}
\begin{minted}[linenos]{asm}
__attribute__((naked)) void SVCall_Handler(void) {
  __asm volatile(
      "tst   lr, #4\n"      // test EXC_RETURN[2]
      "ite   eq\n"          // if-then-else block
      "mrseq r0, msp\n"     // if bit 2 == 0: use MSP
      "mrsne r0, psp\n"     // if bit 2 == 1: use PSP
      "b     SVC_Handler_C");  // jump to C handler
}
\end{minted}

\textbf{Key Implementation Details:}
\begin{itemize}[leftmargin=1.5em]
  \item The \texttt{tst lr, \#4} pattern conforms to ARMv7-M exception semantics: bit~2 of the EXC\_RETURN value identifies whether the interrupted context used MSP (0) or PSP (1).
  \item The \texttt{ite} instruction creates a conditional execution block, allowing both stack selection paths to be handled efficiently.
  \item Passing the stacked frame pointer in r0 to \texttt{SVC\_Handler\_C} avoids memory copying and simplifies return value handling.
  \item The naked attribute prevents the compiler from generating prologue/epilogue code that would interfere with the exception return mechanism.
\end{itemize}

\subsection{C-Level Decoder}
The \texttt{SVC\_Handler\_C} function decodes the SVC instruction and dispatches to the appropriate kernel service.

\vspace{0.3cm}
\noindent\textbf{File:} \file{kern/arch/stm32f446re/sys\_lib/stm32\_startup.c}
\begin{minted}[linenos,fontsize=\footnotesize]{c}
void SVC_Handler_C(uint32_t *stack) {
  // Extract stacked PC to locate the SVC instruction
  uint32_t pc = stack[6];
  
  // Read the SVC instruction opcode (at PC - 2 in Thumb mode)
  uint8_t *svc_inst = (uint8_t *)(pc - 2);
  uint8_t svc_number = svc_inst[0];  // Lower byte contains immediate
  
  // Extract arguments from stacked registers
  uint32_t a0 = stack[0];  // R0
  uint32_t a1 = stack[1];  // R1
  uint32_t a2 = stack[2];  // R2
  uint32_t a3 = stack[3];  // R3
  
  // Dispatch to kernel service
  int32_t retval = syscall_dispatch(svc_number, a0, a1, a2, a3);
  
  // Write return value back to stacked R0
  stack[0] = (uint32_t)retval;
}
\end{minted}

\textbf{Implementation Details:}
\begin{enumerate}[leftmargin=2em,label=\alph*)]
  \item \textbf{SVC Immediate Extraction:} The stacked PC (\texttt{stack[6]}) points to the instruction \emph{after} the SVC. In Thumb-2 mode, an \texttt{svc} instruction is encoded as \texttt{0xDFxx}, where the lower byte (\texttt{xx}) carries the syscall number.\footnote{The ARMv7-M architecture reference specifies this encoding in section A7.7.175.}
  \item \textbf{Argument Extraction:} Hardware automatically stacks \texttt{r0--r3} at positions [0--3], making them directly accessible as function arguments.
  \item \textbf{Return Value Handling:} The return value is written back to \texttt{stack[0]}, which corresponds to register r0. When the exception unwinds, the user code sees this value.
\end{enumerate}

This keeps the handler deterministic and side-effect free---vital when the same path is reused by nested exceptions or future unprivileged tasks.

\section{Syscall Dispatcher}

\subsection{Dispatcher Architecture}
The dispatcher in \file{kern/syscall/syscall.c} serves as the central routing mechanism for all system calls. It maintains a single switch statement with explicit fall-through prevention and ENOSYS handling for invalid syscall numbers.

\begin{figure}[H]
\centering
\begin{tikzpicture}[node distance=1cm, every node/.style={font=\small}]
  \node[process] (entry) {syscall\_dispatch()};
  \node[decision, below=1.2cm of entry, aspect=2.5] (switch) {Syscall\\Number?};
  
  % Left branch
  \node[smallblock, below left=1.5cm and 2.5cm of switch] (read) {SYS\_read\\(50)};
  \node[smallblock, below=0.5cm of read] (write) {SYS\_write\\(55)};
  
  % Center branch
  \node[smallblock, below=1.5cm of switch] (time) {SYS\_time\\(113)};
  \node[smallblock, below=0.5cm of time] (getpid) {SYS\_getpid\\(5)};
  
  % Right branch
  \node[smallblock, below right=1.5cm and 2.5cm of switch] (yield) {SYS\_yield\\(120)};
  \node[smallblock, below=0.5cm of yield] (reboot) {SYS\_reboot\\(119)};
  \node[smallblock, below=0.5cm of reboot] (exit) {SYS\_exit\\(3)};
  
  % Default case
  \node[smallblock, right=2cm of yield, fill=red!20] (enosys) {Return\\-ENOSYS};
  
  % Merge point
  \node[process, below=2.5cm of time] (return) {Return to Handler};
  
  % Arrows
  \draw[line] (entry) -- (switch);
  \draw[line] (switch) -| (read);
  \draw[line] (switch) -| (write);
  \draw[line] (switch) -- (time);
  \draw[line] (switch) -- (getpid);
  \draw[line] (switch) -| (yield);
  \draw[line] (switch) -| (reboot);
  \draw[line] (switch) -| (exit);
  \draw[line] (switch) -- (enosys);
  
  \draw[line] (read) |- (return);
  \draw[line] (write) |- (return);
  \draw[line] (time) -- (return);
  \draw[line] (getpid) -- (return);
  \draw[line] (yield) |- (return);
  \draw[line] (reboot) |- (return);
  \draw[line] (exit) |- (return);
  \draw[line] (enosys) |- (return);
\end{tikzpicture}
\caption{Syscall Dispatcher Flow}
\label{fig:dispatcher}
\end{figure}

The file also hosts \texttt{g\_current\_task\_id}, a global variable exported for both kernel book-keeping and the \texttt{SYS\_getpid} result. This variable is initialized to 1000 during system boot.

\subsection{Implemented Services}
Table~\ref{tab:syscalls} documents each implemented syscall end-to-end. The `User Wrapper` column references the API in \file{userland/utils/unistd.c}; the kernel actions cite the core functions touched.

\begin{table}[h]
  \centering
  \begin{tabularx}{\linewidth}{@{}L{2.7cm}C{1.2cm}Y Y@{}}
    \toprule
    \textbf{Syscall} & \textbf{ID} & \textbf{User Wrapper} & \textbf{Kernel Action} \\
    \midrule
    \texttt{SYS\_read} & 50 & \texttt{ssize\_t read(int, void*, size)} & Block on console RX (\texttt{IsDataAvailable}) until newline or 256 bytes, terminate buffer, return byte count. \\
    \texttt{SYS\_write} & 55 & \texttt{ssize\_t write(int, const void*, size)} & Iterate through buffer, transmit via \texttt{Uart\_write} for \texttt{STDOUT\_FILENO}. \\
    \texttt{SYS\_\_\_time} & 113 & \texttt{uint32\_t getSysTickTime(void)} & Return \texttt{\_\_getTime()} (ms counter). \\
    \texttt{SYS\_getpid} & 5 & \texttt{int getpid(void)} & Expose \texttt{g\_current\_task\_id} (seeded to 1000). \\
    \texttt{SYS\_yield} & 120 & \texttt{int yield(void)} & Set \reg{SCB->ICSR |= PENDSVSET}; return 0. \\
    \texttt{SYS\_reboot} & 119 & \texttt{int reboot(void)} & Invoke \texttt{\_\_NVIC\_SystemReset()} for a full MCU reset. \\
    \texttt{SYS\_\_exit} & 3 & \texttt{void exit(int)} & Pend PendSV for scheduler to reclaim task; stub loops in \texttt{\_exit}. \\
    \bottomrule
  \end{tabularx}
  \caption{Implemented syscalls and their execution path.}
  \label{tab:syscalls}
\end{table}

Unknown call numbers fall back to \texttt{ENOSYS}, enabling straightforward validation from userland.

\subsection{Support Utilities}
\begin{itemize}[leftmargin=1.5em]
  \item UART services rely on the driver trio \file{kern/arch/stm32f446re/sys\_lib/sys\_usart.c}, \file{kern/lib/UsartRingBuffer.c}, and \file{kern/lib/kern/serial\_lin.c} for buffered I/O.
  \item Timing relies solely on the SysTick counter---no additional hardware timers are consumed.
  \item \texttt{g\_current\_task\_id} lives in RAM (BSS); a volatile pointer alias (\texttt{\_\_g\_current\_task\_id\_ref}) exists to prevent dead-code elimination and simplify debugger inspection.
\end{itemize}

\section{Userland Interface}

\subsection{Wrapper Macro}
All syscall wrappers share a common \texttt{SVC\_CALL} macro that provides a clean interface between user code and the kernel. This macro uses inline assembly to comply with ARM EABI calling conventions.

\vspace{0.3cm}
\noindent\textbf{File:} \file{userland/utils/unistd.c}
\begin{minted}[linenos,fontsize=\footnotesize]{c}
#define SVC_CALL(num, a0, a1, a2, a3) ({               \
  register uint32_t r0 __asm("r0") = (uint32_t)(a0);    \
  register uint32_t r1 __asm("r1") = (uint32_t)(a1);    \
  register uint32_t r2 __asm("r2") = (uint32_t)(a2);    \
  register uint32_t r3 __asm("r3") = (uint32_t)(a3);    \
  __asm volatile("svc %[imm]"                           \
                 : "+r"(r0)                             \
                 : [imm] "I"(num), "r"(r1), "r"(r2), "r"(r3) \
                 : "memory");                           \
  r0; })
\end{minted}

\textbf{Key Properties:}
\begin{itemize}[leftmargin=1.5em]
  \item \textbf{Register Constraints:} The \texttt{\_\_asm("r0")} syntax forces the compiler to place arguments in specific registers (r0--r3), matching the ARM EABI.
  \item \textbf{Return Value Capture:} The inline assembly constraint \texttt{"+r"} for \reg{r0} indicates it is both an input and output, capturing the kernel return value without extra moves.
  \item \textbf{Memory Barrier:} The explicit clobber of \texttt{"memory"} prevents the compiler from reordering memory accesses around the syscall, ensuring proper sequencing of operations.
  \item \textbf{Immediate Constraint:} The \texttt{"I"} constraint requires the syscall number to be a compile-time constant, allowing it to be encoded directly in the SVC instruction.
  \item \textbf{EABI Compliance:} All wrappers comply with the standard ARM EABI argument passing rules, ensuring compatibility with any future compiled user applications.
\end{itemize}

\subsection{Example Wrapper Implementation}
Here's the implementation of the \texttt{write()} syscall wrapper:

\begin{minted}[linenos,fontsize=\footnotesize]{c}
ssize_t write(int fd, const void *buf, size_t count) {
  return (ssize_t)SVC_CALL(SYS_write, fd, buf, count, 0);
}
\end{minted}

The wrapper is simple and efficient---it merely packages the arguments and invokes the SVC instruction. All complexity is handled by the kernel dispatcher.

\subsection{Published API}
The user-visible prototypes reside in \file{userland/include/unistd.h}. Besides the mandatory seven wrappers, \texttt{\_exit()} spins after pended termination to guarantee that the caller never returns to a potentially reclaimed stack frame.

\begin{minted}[linenos,fontsize=\footnotesize]{c}
// userland/include/unistd.h
ssize_t read(int fd, void *buf, size_t count);
ssize_t write(int fd, const void *buf, size_t count);
uint32_t getSysTickTime(void);
int getpid(void);
int yield(void);
int reboot(void);
void exit(int status) __attribute__((noreturn));
\end{minted}

\section{Demo Application}

\subsection{Test Application Overview}
The demonstrator in \file{kern/kmain/kmain.c} performs an A-to-Z validation of the syscall surface. The application systematically exercises each implemented syscall to verify correct operation.

\begin{figure}[H]
\centering
\begin{tikzpicture}[node distance=0.9cm, every node/.style={font=\small}]
  \node[cloud, text width=4cm] (boot) {System Boot\\\_\_sys\_init()};
  \node[process, below=of boot, text width=5cm] (init) {Initialize task\_id = 1000};
  \node[process, below=of init, text width=5cm] (write) {Test write()\\Print greeting to UART};
  \node[process, below=of write, text width=5cm] (time) {Test getSysTickTime()\\Query system timer};
  \node[process, below=of time, text width=5cm] (getpid) {Test getpid()\\Get current task ID};
  \node[process, below=of getpid, text width=5cm] (yield) {Test yield()\\Trigger PendSV};
  \node[process, below=of yield, text width=5cm, fill=red!10] (exit) {Test exit(0)\\Terminate task};
  
  \draw[line] (boot) -- (init);
  \draw[line] (init) -- (write);
  \draw[line] (write) -- (time);
  \draw[line] (time) -- (getpid);
  \draw[line] (getpid) -- (yield);
  \draw[line] (yield) -- (exit);
  
  % Side annotations
  \node[right=0.5cm of write, text width=3.5cm, align=left, font=\footnotesize] {Verifies UART\\output path};
  \node[right=0.5cm of time, text width=3.5cm, align=left, font=\footnotesize] {Confirms timing\\infrastructure};
  \node[right=0.5cm of getpid, text width=3.5cm, align=left, font=\footnotesize] {Tests global\\state access};
\end{tikzpicture}
\caption{Demo Application Execution Flow}
\label{fig:demo-flow}
\end{figure}

\subsection{Test Sequence}
\begin{enumerate}[leftmargin=2em,label=\arabic*.]
  \item \textbf{Hardware Initialization:} Call \texttt{\_\_sys\_init()} and assign \texttt{task\_id = 1000}.
  \item \textbf{UART Output Test:} Print a greeting via \texttt{write(STDOUT, ...)} and confirm UART TX.
  \item \textbf{Timing Test:} Query \texttt{getSysTickTime()} and verify the millisecond counter is incrementing.
  \item \textbf{Process ID Test:} Call \texttt{getpid()} and verify it returns 1000; print both time and PID via kernel \texttt{kprintf()} for visibility on the debug UART.
  \item \textbf{Scheduling Test:} Trigger \texttt{yield()} to demonstrate that PendSV requests are raised from both userland and SysTick (every 10~ms).
  \item \textbf{Termination Test:} Finish with \texttt{exit(0)}, after which control never returns because PendSV will eventually recycle the context once the scheduler is in place.
\end{enumerate}

Throughout this flow, the kernel prints diagnostic lines confirming each stage, making it trivial to capture UART logs as evidence of correct operation.

\section{Memory Usage Analysis}

\subsection{Linker Statistics}
The following memory statistics were extracted from the linker map file (\texttt{compile/Makefile} output) for the current build:

\begin{table}[h]
\centering
\begin{tabular}{lrrr}
\toprule
\textbf{Section} & \textbf{Size (bytes)} & \textbf{Size (KB)} & \textbf{Percentage} \\
\midrule
Text (Flash) & 20,948 & 20.5 & 4.0\% of 512 KB \\
Data (RAM init) & 108 & 0.1 & 0.08\% of 128 KB \\
BSS (RAM zero) & 2,305 & 2.3 & 1.8\% of 128 KB \\
\midrule
\textbf{Total Flash} & \textbf{20,948} & \textbf{20.5} & \textbf{4.0\%} \\
\textbf{Total RAM} & \textbf{2,413} & \textbf{2.4} & \textbf{1.9\%} \\
\bottomrule
\end{tabular}
\caption{Memory utilization breakdown}
\label{tab:memory}
\end{table}

\subsection{Analysis}
These numbers were collected from the map file emitted via \texttt{compile/Makefile}. The small footprint demonstrates efficient implementation:

\begin{itemize}[leftmargin=1.5em]
  \item \textbf{Flash Usage:} Only 4.0\% of available flash is consumed, leaving 491.5 KB (96\%) available for the upcoming scheduler, IPC primitives, and filesystem services.
  \item \textbf{RAM Usage:} Minimal RAM consumption (1.9\%) ensures ample space for task stacks, dynamic memory allocation, and kernel data structures.
  \item \textbf{Code Density:} The implementation achieves comprehensive syscall functionality with minimal overhead, validating the design choices.
\end{itemize}

\section{Testing and Verification}

\subsection{Test Methodology}
All syscalls were tested individually and in integrated scenarios using both hardware debugging and serial terminal monitoring. The test platform consisted of the STM32F446RE Nucleo board connected via ST-LINK debugger and UART to a host PC.

\subsection{Test Results}

\begin{table}[h]
\centering
\begin{tabularx}{\linewidth}{@{}l X c@{}}
\toprule
\textbf{Test Category} & \textbf{Description} & \textbf{Status} \\
\midrule
UART I/O & Verified bidirectional character flow for \texttt{read()} and \texttt{write()} using a serial terminal at 115200 baud; newline termination and buffer capping (256 bytes) behave as designed. & \textcolor{green!60!black}{\textbf{PASS}} \\[0.3cm]
Timing Infrastructure & Observed the millisecond counter (\texttt{g\_msTick}) incrementing monotonically and the 10~ms PendSV cadence on a debugger watch window. Verified timing accuracy with oscilloscope. & \textcolor{green!60!black}{\textbf{PASS}} \\[0.3cm]
System Reset & Invoked \texttt{reboot()} from userland; the MCU executed NVIC system reset and returned to \texttt{kmain()} cleanly with all globals and peripherals reinitialised. & \textcolor{green!60!black}{\textbf{PASS}} \\[0.3cm]
Task Management & Confirmed that \texttt{yield()} triggers the PendSV stub (verified via debugger breakpoint) and that \texttt{exit()} never returns to the caller, matching a cooperative multitasking contract. & \textcolor{green!60!black}{\textbf{PASS}} \\[0.3cm]
Process ID & Verified \texttt{getpid()} consistently returns the seeded task ID (1000) across multiple invocations. & \textcolor{green!60!black}{\textbf{PASS}} \\[0.3cm]
Error Handling & Tested invalid syscall numbers; confirmed dispatcher returns \texttt{-ENOSYS} (38) as expected. & \textcolor{green!60!black}{\textbf{PASS}} \\
\bottomrule
\end{tabularx}
\caption{Comprehensive test results}
\label{tab:test-results}
\end{table}


\section{Known Limitations and Future Work}

\subsection{Current Limitations}
\begin{enumerate}[leftmargin=2em]
  \item \textbf{Scheduler Stub:} PendSV currently serves as a stub that immediately returns. The context switcher and Task Control Block (TCB) management must be implemented to enable true multitasking.
  
  \item \textbf{Limited File Descriptors:} The syscall layer presently filters solely on \texttt{STDIN\_FILENO}/\texttt{STDOUT\_FILENO}. Extending to additional devices or descriptors will require a dispatch table and resource tracking infrastructure.
  
  \item \textbf{Minimal Error Handling:} Error codes are basic. Future work should propagate comprehensive errno values (\texttt{-EINTR}, \texttt{-EBADF}, \texttt{-EAGAIN}, etc.) and integrate with a userspace-visible \texttt{errno} mechanism.
  
  \item \textbf{No Privilege Separation:} The demo executes in privileged mode. Once multiple tasks exist, PSP/MSP segregation and unprivileged thread mode will be introduced to enforce proper isolation.
  
  \item \textbf{Blocking I/O:} The \texttt{read()} syscall blocks indefinitely. A timeout mechanism or non-blocking variants would improve responsiveness.
  
  \item \textbf{Single Task Model:} Currently supports only one implicit task. Multi-task support requires scheduler, stack management, and inter-task communication primitives.
\end{enumerate}

\section{Implementation File Inventory}

\subsection{Core System Call Infrastructure}
\begin{table}[h]
\centering
\begin{tabularx}{\linewidth}{@{}l X@{}}
\toprule
\textbf{File} & \textbf{Description} \\
\midrule
\file{kern/syscall/syscall.c} & Central dispatcher implementing switch-based routing for all seven syscalls; hosts \texttt{g\_current\_task\_id} \\[0.2cm]
\file{kern/include/syscall.h} & Syscall number definitions, dispatcher prototypes, and kernel-side type declarations \\[0.2cm]
\file{userland/utils/unistd.c} & User-space wrapper implementations using SVC\_CALL macro for all seven syscalls \\[0.2cm]
\file{userland/include/unistd.h} & Public API declarations for user applications; POSIX-style function signatures \\
\bottomrule
\end{tabularx}
\end{table}

\subsection{Exception Handling and Hardware Integration}
\begin{table}[h]
\centering
\begin{tabularx}{\linewidth}{@{}l X@{}}
\toprule
\textbf{File} & \textbf{Description} \\
\midrule
\file{kern/arch/stm32f446re/sys\_lib/stm32\_startup.c} & Vector table, naked SVCall\_Handler (assembly shim), and SVC\_Handler\_C (decoder) \\[0.2cm]
\file{kern/arch/cm4/systick.c} & SysTick ISR, millisecond counter maintenance, 10~ms PendSV trigger logic \\[0.2cm]
\file{kern/lib/kern/sys\_init.c} & System initialization: clock configuration, FPU enable, UART setup, SysTick initialization \\
\bottomrule
\end{tabularx}
\end{table}

\subsection{Support Libraries and Drivers}
\begin{table}[h]
\centering
\begin{tabularx}{\linewidth}{@{}l X@{}}
\toprule
\textbf{File} & \textbf{Description} \\
\midrule
\file{kern/arch/stm32f446re/sys\_lib/sys\_usart.c} & Low-level UART driver for STM32F446RE (USART2 \& USART6) \\[0.2cm]
\file{kern/lib/UsartRingBuffer.c} & Circular buffer implementation for UART RX/TX buffering \\[0.2cm]
\file{kern/lib/kern/serial\_lin.c} & Serial line discipline, character echo, line editing \\[0.2cm]
\file{kern/include/kern/types.h} & Type definitions and constants (errno values, file descriptors) \\[0.2cm]
\file{kern/kmain/kmain.c} & Demo application orchestrating syscall validation tests \\
\bottomrule
\end{tabularx}
\end{table}

\subsection{Build System}
\begin{itemize}[leftmargin=1.5em]
  \item \file{compile/Makefile} -- Master build script with compiler flags, linker script integration, and map file generation
\end{itemize}

\section{Conclusion}

This assignment successfully establishes a robust, hardware-accelerated system call infrastructure for DUOS on the STM32F446RE platform. The implementation demonstrates several key achievements:


\vspace{1cm}
\begin{center}
\rule{0.8\linewidth}{0.5pt}\\[0.3cm]
\textit{End of Report}
\end{center}

\end{document}
